% File: Volume-II-Applications/appendices/appH-observational-data.tex

\chapter{Observational Data Sources}
\label{app:obs-data}

This appendix summarizes the principal observational data sources used
throughout Volume~II.  Our aim is not to provide an exhaustive review
of the observational literature, but to give sufficient detail for
interpreting the quantitative comparisons performed in
Chapters~8--10.  We organize the discussion by data type: cosmic
microwave background (CMB), large--scale structure (LSS), and
supernovae (SNe).  Additional references are provided at the end of
the chapter.

\section{Cosmic Microwave Background}

For the CMB temperature and polarization anisotropies, the principal
data sets come from \emph{Planck 2018}, including power spectra
$C_\ell^{TT}$, $C_\ell^{TE}$, and $C_\ell^{EE}$.  These are obtained
from full--sky maps after foreground subtraction and masking.  Noise
properties, beam transfer functions, and sky cuts are supplied by the
Planck Legacy Archive.

In \rsvpabbr{}, the primary modifications appear through:
\begin{itemize}
  \item the lamphrodynamic evolution of scalar perturbations,
  \item entropy--vector couplings affecting anisotropic stress,
  \item modified late--time integrated Sachs--Wolfe contributions.
\end{itemize}
High--$\ell$ damping tails are sensitive to recombination physics, but
the dominant lamphrodynamic signatures enter through low--$\ell$
multipoles and potential deviations in the acoustic peaks.

\section{Large--Scale Structure and BAO}

For LSS, we rely on galaxy surveys including BOSS, eBOSS, DESI, and
future surveys such as Euclid and LSST.  The primary observables are
the matter power spectrum $P(k)$, baryon acoustic oscillations (BAO),
and weak lensing shear correlations.  

Lamphrodynamic modifications arise from:
\begin{itemize}
  \item altered linear growth factors,
  \item modified transfer functions for density perturbations,
  \item possible effective ``dark--matter--like'' behavior from
    entropy--vector dynamics.
\end{itemize}
Comparison with $\Lambda$CDM predictions requires careful statistical
analysis (see Appendix~G) and numerical integration (Appendix~I).

\section{Galaxy Clusters and Weak Lensing}

Galaxy clusters probe the nonlinear regime of structure formation.
Mass estimates rely on X--ray temperatures, Sunyaev--Zel'dovich
signals, and weak lensing measurements.  In \rsvpabbr{}, cluster
profiles provide sensitive tests of emergent--dark--matter behavior and
potential deviations in the mass--concentration relation.  Weak
lensing convergence maps probe projected mass densities and are
significantly influenced by modified geodesic equations (Chapter~3).

\section{Supernovae and Standard Candles}

Type~Ia supernovae provide direct constraints on the expansion
history through distance moduli $\mu(z)$.  The most comprehensive
compilations are the Pantheon and Pantheon+ samples.  In \rsvpabbr{},
the relation between redshift and luminosity distance is modified by
entropy--induced redshifting mechanisms (Chapter~4).  Systematic
uncertainties include dust extinction, selection effects, and
population drift.

\section{Cosmic Chronometers}

Cosmic chronometers measure the expansion rate directly from relative
ages of passively evolving galaxies.  This approach provides
model--independent constraints on $H(z)$ and is therefore particularly
useful for models with modified redshift relations.  The current
precision is limited by stellar population modeling and metallicity
effects.

\section{Future Surveys}

Future surveys (Euclid, LSST, \emph{Simons Observatory}, CMB--S4,
LIGO--Voyager) will dramatically extend the volume and precision of
observations.  These are of particular relevance to \rsvpabbr{} due to:
\begin{itemize}
  \item increased sensitivity to large--scale entropy modes,
  \item improved constraints on growth rates and anisotropic stress,
  \item ability to discriminate lamphrodynamic perturbations from
    modified gravity and dark sector models.
\end{itemize}

\section{Data Availability}

Most observational data sets are publicly available:
\begin{itemize}
  \item Planck Legacy Archive (CMB),
  \item SDSS/BOSS/eBOSS data releases (galaxy surveys),
  \item DESI data products (spectroscopic survey),
  \item Pantheon/Pantheon+ (supernovae),
  \item public weak lensing catalogs (DES, KiDS).
\end{itemize}
Where proprietary restrictions exist, we provide references and key
summaries suitable for comparative analysis.

\section{Remarks and Further Reading}

For detailed discussions of observational data, see the Planck 2018
papers, DESI survey documentation, and the Pantheon+ release papers.
For weak lensing and cluster analyses, see recent DES, KiDS, and
X--ray/SZ collaborations.  Volume~II uses these sources as benchmarks
for the observational viability of the lamphrodynamic dynamics of
\rsvpabbr{}.

